\section{Robotics、AGI 时间表与经济影响：技术乐观与现实摩擦并存}

\subsection{Embodied AI 的高约束世界}

对话对 robotics 的判断是“有机会，但门槛极高”。在数字系统里，失败可快速回滚；在物理系统里，错误会变成安全事故。Lex 用一句话概括了这个约束：\textit{``embodied systems are almost allowed to fail never''}。这让机器人落地节奏明显慢于纯软件 Agent。

\begin{importantbox}{为什么自动驾驶和工业自动化更先落地}
\begin{itemize}[leftmargin=1.5em]
    \item 任务边界更清晰，评价指标可定义；
    \item 运行环境相对可控，可通过基础设施降低不确定性；
    \item 经济回报可量化，投入与收益路径更明确。
\end{itemize}
\end{importantbox}

\subsection{AGI 时间表：分歧的根源在定义}

访谈没有给出统一时间表，核心原因是 AGI 定义不一致。若定义为“完成大部分数字经济任务”，时间可能更近；若定义为“在所有认知维度稳定超人”，时间明显更远。Nathan 反复强调“jagged intelligence”特征：模型在某些任务超强，在另一些任务很脆弱。

\begin{warningbox}{“超人程序员”标签可能掩盖系统性缺口}
把少数 benchmark 高分等同于“可替代复杂工程角色”会产生误判。真实软件生产包含组织协作、需求博弈、跨团队沟通、长期维护责任，这些都不是单轮代码生成可以覆盖的。
\end{warningbox}

\subsection{本章小结}

AGI 讨论在 2026 仍应保持工程化态度：先看具体任务闭环，再谈宏大标签。机器人与全面自动化不会线性到来，而会按场景逐步渗透。

